\newpage
\section{Marco legal}
\subsection{Ley Universitaria Nro. 30220}
La Ley Universitaria N° 30220, promulgada el 9 de julio de 2014, establece en su Artículo 49 que las universidades deben promover las Prácticas Pre - Profesionales de sus estudiantes con el objetivo de vincular la formación académica con la actividad laboral y productiva del país. El Artículo 90 dispone que los programas de estudio incluyan periodos de prácticas que fortalezcan la formación integral del egresado y faciliten su inserción al mercado laboral mediante el desarrollo de competencias aplicadas.

\subsection{Estatuto de la Universidad Nacional del Altiplano}
El Estatuto de la Universidad Nacional del Altiplano Puno, en concordancia con la Ley Universitaria N° 30220, establece que la formación profesional de los estudiantes debe articular los conocimientos teóricos con la experiencia práctica. Dispone que las Escuelas Profesionales organicen y supervisen las prácticas Pre - Profesionales como requisito para la obtención del grado de bachiller, garantizando el acompañamiento docente y que las actividades correspondan al perfil de la carrera profesional.

\subsection{Reglamento de Prácticas Pre - Profesionales}
El Reglamento de Prácticas Pre - Profesionales de la Escuela Profesional de Ingeniería de Sistemas (Código RG-001-FIMEES-P29-R, versión V01, vigente desde el 05 de agosto de 2024) regula el proceso de inscripción, desarrollo, supervisión y evaluación de las prácticas. Establece una duración mínima de cuatro meses, define los roles del practicante, el jefe del área receptora y el docente supervisor, y fija los instrumentos oficiales de evaluación (actitudinal, desempeño y rúbrica de sustentación) que deben completarse para la aprobación del informe final.
